\chapter{Этаж~0. Планковский пол}
\label{ch:floor0}
\label{ch:floor0-spectrum}

\section{Один планкон в одной ячейке}
\label{sec:floor0-object}

Этаж~0: один планкон в одной планковской ячейке~$\hv=\PlanckCell$
(\cref{ch:floors-preface}).
Занятость $b=1$, топологический заряд $n=\pm1$ на контуре (\cref{sec:planckon-floor0,ch:matter}).
Это не картина «орбитали внутри радиуса $r<\ell_P$'' --- ниже шага~$\caStepSpace$ пространства как континуума нет.
И не этаж~2 (кварки, составные квантовые числа): здесь только внутренние степени свободы
\emph{одного} локализованного возбуждения.

На макроуровне~$\layerMacro$ атомные уровни и орбитали --- то, как прибор считывает составной узел;
на~$\layerMicro$ у планкона заранее действует та же \emph{зонная логика},
что у атома, ядра или кристалла: дискретные уровни, правила отбора, запрет Паули,
конечный набор состояний (бит-бюджет $\caBitBudget$, \cref{ch:evolution}).

\begin{table}[htbp]
  \centering
  \caption{Аналогия: атом / ядро на~$\layerMacro$ и планкон на~$\layerMicro$ (этаж~0).}
  \label{tab:floor0-analogy}
  \begin{tabular}{@{}ll@{}}
    \toprule
    Атом, ядро ($\layerMacro$) & Планкон ($\layerMicro$) \\
    \midrule
    уровни $E_n$ & $E=n_E \ladderEnergy$, $n_E\in\mathbb{Z}$; фаза $\Phi$ на кольце \\
    орбитали $(n,\ell,m)$ & внутренние: кольцо + $SU(2)$; внешние: оболочки соседей (\cref{ch:floor1}) \\
    запрет Паули & A16: два совпадающих спинора в одном $\hv$ запрещены; $b\in\{0,1\}$ \\
    запрещённый переход & скачок не кратен $\ladderEnergy$, $\Delta\varphi_{\min}$ или ломает $n$ без партнёра \\
    запрещённая зона (кристалл) & щели $\omega(k)$ на зоне Бриллюэна вакуума (\cref{ch:matter}) \\
    \bottomrule
  \end{tabular}
\end{table}

\section{Два внутренних регистра}
\label{sec:floor0-registers}

На одной ячейке независимых «полей Стандартной модели» нет.
Из аксиом и закона~$g$ (\cref{ch:axioms,ch:evolution}) следуют два связанных регистра:

\paragraph{Фазовое кольцо.}
Накопленная фаза импульсного реестра
$\Phi\in\mathbb{Z}_{\phaseRingCount}$, $\phaseRingCount=512$.
Один полный оборот $U(1)$ на кольце --- $512$ дискретных тиков.
Минимальный разрешённый скачок $\Delta\varphi_{\min}=\tfrac12$ рад.
Из леммы~\ref{lem:51a} следует
\[
  \Delta\varphi_{\mathrm{disc}}
  = \left\lfloor \phaseRingCount\,\frac{\Delta\varphi_{\min}}{2\pi}\right\rfloor
  = 41
\]
тик на $\mathbb{Z}_{\phaseRingCount}$.
Один такой скачок --- один квант энергии $\ladderEnergy$ лестницы (\cref{thm:51,lem:51e}).

\paragraph{Спинор на узле планкона.}
Состояние $z\in\mathbb{C}^2$ в ячейке~$\hv$, где планкон ($b=1$).
Вращение вокруг оси Блоха на этом узле реализует $SU(2)$:
при повороте на $2\pi$ состояние меняет знак, на $4\pi$ --- возвращается
(полуцелый спин, A16).
Это \emph{внутренняя} ориентация, не топологический заряд $n$ на контуре.

Оба регистра живут на одном носителе и ограничены одним бит-бюджетом;
их нельзя трактовать как два произвольных квантовых числа лабораторной таблицы.

\section{Лестница $n_E$ на кольце}
\label{sec:floor0-ne-ladder}

Энергия события на~$\layerMicro$ квантуется лестницей $E=n_E \ladderEnergy$ (\cref{lem:51e}).
Число $n_E$ считается по накопленным тикам фазы на кольце:
\[
  n_E = \left\lfloor \frac{|\Phi|}{\Delta\varphi_{\mathrm{disc}}}\right\rfloor,
  \qquad
  \Delta\varphi_{\mathrm{disc}} = 41.
\]

\begin{table}[htbp]
  \centering
  \caption{Вехи внутреннего спектра на $\mathbb{Z}_{\phaseRingCount}$.}
  \label{tab:floor0-landmarks}
  \begin{tabular}{@{}rlcl@{}}
    \toprule
    веха & тики $\Phi$ & $n_E$ & смысл \\
    \midrule
    первый квант & $41$ & $1$ & минимальный разрешённый скачок ($\Delta\varphi_{\min}$) \\
    второй квант & $82$ & $2$ & второй уровень лестницы \\
    полуоборот Паули & $256$ & $6$ & импульс $\phaseRingCount/2$ (поворот на $\pi$) \\
    полный оборот кольца & $512$ & $12$ & цикл $U(1)$ на $\mathbb{Z}_{\phaseRingCount}$ \\
    \bottomrule
  \end{tabular}
\end{table}

Связь с запретом Паули: импульс $\pi$ на половине кольца ($256$ тиков) --- та же дискретная
$\pi$-ротация, что в конгруэнтной лестнице CL-2 (\cref{ch:evolution}).
Полный оборот $512$ тиков даёт $n_E=12$ --- двенадцать квантов $\ladderEnergy$ на один цикл фазы.

\begin{proposition}[Лестница $n_E$]
\label{prop:floor0-ne-ladder}
Для $\phaseRingCount=512$ и $\Delta\varphi_{\mathrm{disc}}=41$:
\[
  n_E(\Phi) = \lfloor |\Phi| / 41\rfloor,\quad
  \Phi_{\pi} = \phaseRingCount/2 = 256,\quad
  \Phi_{2\pi} = \phaseRingCount = 512.
\]
\end{proposition}

\begin{proof}
\pfrom{определения $n_E$ и $\Delta\varphi_{\mathrm{disc}}=41$}
первые две строки таблицы~\ref{tab:floor0-landmarks} дают $n_E=1$ при $|\Phi|=41$ и $n_E=2$ при $|\Phi|=82$.
\pusing{CL-2 (\cref{ch:evolution}): половина кольца --- поворот на $\pi$}
$|\Phi|=\phaseRingCount/2=256$ и $256/41=6$ с остатком, откуда $n_E=6$.
\pusing{полного оборота $|\Phi|=\phaseRingCount$}
$512/41=12$ с остатком, откуда $n_E=12$ и $\Phi_{2\pi}=512$.
\end{proof}

\section{Основное состояние}
\label{sec:floor0-ground}

Устойчивый планкон с $n=\pm1$ после релаксации поля (вихрь на одном~$\hv$)
имеет следующие свойства на ядре и контуре:

\begin{itemize}[noitemsep]
\item $b=1$ --- ячейка занята планконом, не планковской дыркой;
\item на контуре устойчивый вихревой паттерн с $|n|$ не ниже порога автоматической оценки топозаряда;
\item $|\Delta\varphi|_{\mathrm{core}} < \Delta\varphi_{\min}$ --- фазовый дефект сидит
  в дне Гейзенберга, без острого скачка аргумента;
\item $n_E=0$ в реестре энергии: за время релаксации не накопилось
  $\ge 41$ тиков насыщающей фазы на ядре;
\item вектор Блоха на ядре устойчив под повторным применением~$g$
  (внутренняя ось $SU(2)$ не дрейфует).
\end{itemize}

\begin{proposition}[Основное состояние планкона]
\label{prop:floor0-vortex-ground}
После релаксации вихря с $n=\pm1$ на одном~$\hv$:
$b=1$, $|\Delta\varphi|_{\mathrm{core}}<\Delta\varphi_{\min}$, $n_E=0$,
устойчивый вектор Блоха на ядре.
\end{proposition}

\begin{proof}
\pfrom{леммы~\ref{lem:51a} и схемы импульсного реестра (\cref{ch:evolution})}
скачок фазы меньше $\Delta\varphi_{\min}$ на ядре не даёт целого кванта $\ladderEnergy$,
поэтому $n_E=\lfloor|\Phi|/41\rfloor=0$ после релаксации.
\pfrom{локализации планкона (\cref{sec:planckon-floor0})}
$b=1$ на ядре.
\pfrom{устойчивости вихря при $n=\pm1$}
контур несёт топологический заряд; ядро остаётся в допустимой зоне ротора $SU(2)$.
\end{proof}

Это \emph{нулевой} внутренний уровень по $n_E$ при ненулевом топологическом $n$:
планкон уже локализован и массивен (\cref{ch:matter}), но фазовый реестр на кольце
в основном состоянии пуст.

\section{Монодромия \texorpdfstring{$SU(2)$}{SU(2)} на ядре}
\label{sec:floor0-su2}

Топологический $n$ отвечает за обход контура; спин --- за поведение спинора \emph{на ядре}.
Вращение $R_\alpha$ вокруг оси Блоха ядра даёт скалярное произведение
\[
  \langle z,\, R_{2\pi} z\rangle \approx -1,
  \qquad
  \langle z,\, R_{4\pi} z\rangle \approx +1.
\]
Первое --- знак минус при одном обороте (фермионная монодромия);
второе --- возврат после двух оборотов.
Это тот же смысл, что у спина~$\tfrac12$ в квантовой механике,
но здесь проверяется на дискретном спиноре~$z$ в центре~$\hv$, а не постулируется.

\begin{proposition}[Полуцелый спин на ядре]
\label{prop:floor0-su2-monodromy}
Для устойчивого ядра вихря: $2\pi\mapsto -1$, $4\pi\mapsto +1$ в смысле
$\mathrm{Re}\langle z, R_\phi z\rangle$.
\end{proposition}

\begin{proof}
\pfrom{CL-2 (\cref{ch:evolution}): $R(N_{\mathrm{ring}}/2)=-\mathbb{1}$ на спиноре}
поворот на $\pi$ в дискретном роторе меняет знак $z$.
\pfrom{двойного оборота $4\pi$}
композиция двух таких полуоборотов возвращает знак, откуда $\langle z,R_{4\pi}z\rangle\approx+1$.
\end{proof}

\section{Сколько состояний у одного планкона}
\label{sec:floor0-catalog}

Три внутренних регистра (\cref{sec:floor0-registers}) можно перечислить по отдельности:

\begin{itemize}[noitemsep]
\item \textbf{кольцо} --- $512$ положений $\varphi_{\mathrm{disc}}$ на $\mathbb{Z}_{\phaseRingCount}$;
  канонически это пара $(k_\varphi,\varphi_f)$ с $k_\varphi\in\mathbb{Z}_{13}$,
  $\varphi_f\in\mathbb{Z}_{41}$ и швом в $21$ тик
  ($21=N_\varphi+N_{\mathrm{hier}}$, \cref{ch:evolution});
\item \textbf{лестница $n_E$} --- $13$ классов $0\ldots 12$
  ($n_E=\lfloor|\Phi|/41\rfloor$, полный оборот кольца даёт $n_E=12$);
\item \textbf{ориентация $SU(2)$} --- точка на сфере Блоха;
  на дискретной сетке амплитуд с шагом $2^{-6}$ различных направлений порядка $3\cdot 10^6$.
\end{itemize}

Наивное произведение этих чисел даёт миллиарды комбинаций.
Однако вся ячейка $\hv$ ограничена бит-бюджетом Бекенштейна
$2^{\caBitBudget}\approx 535$ (\cref{ch:evolution}): столько различимых состояний
\emph{всего}, а не на каждый регистр отдельно.
Регистры связаны через представление $Z\in\mathbb{Z}_{\phaseRingCount}[i]$, а не перемножаются
как независимые квантовые числа лабораторной таблицы.

При эволюции закона~$g$ на кипящем вакууме (A5) перечислены орбиты внутренней пары $(q,p)$
на одном узле: различаются основная ветка ($n_E=0$) и возбуждённые ветки ($n_E\ge1$).
Полный перебор всех $2^{B_{\hv}}$ классов по фазе и Блоху для этажа~0 не требуется.

\section{Фазовое пространство одного планкона}
\label{sec:floor0-phase-space}

В модели уже есть обобщённые координаты $\caCoordsDhDt$ и импульсный квант $\ladderMomentum=\ladderAction/\caStepSpace$
(\cref{ch:evolution}).
На этаже~0, когда внешняя позиция $x$ зафиксирована, остаётся дискретное
\textbf{внутреннее} пространство \emph{планкона} $\Gamma_{\hv}$ (бюджет ячейки~$\hv$ на узле планкона) --- пара $(q,p)$:

\begin{itemize}[noitemsep]
\item \textbf{координаты $q$:} положение на кольце $\varphi_{\mathrm{disc}}\in\mathbb{Z}_{\phaseRingCount}$,
  каноническая пара $(k_\varphi,\varphi_f)$, ориентация спинора на сфере Блоха, плотность $|z|^2$;
\item \textbf{импульс $p$:} импульс $\Phi$ за один такт (тики на $\mathbb{Z}_{\phaseRingCount}$, с полом Гейзенберга),
  число $n_E$ в реестре энергии, дискретный $\symPi/\ladderMomentum$ на узле планкона.
\end{itemize}

Наивное произведение размерностей $q$ и $p$ на порядки больше потолка
$2^{\caBitBudget}\approx 535$: $\Gamma_{\hv}$ \emph{сжато} бит-бюджетом ячейки,
а не является декартовым произведением независимых квантовых чисел.

\paragraph{Среда: кипящий океан.}
Решётка без пропусков узлов: на каждой~$\hv$ есть состояние (обычно планковская дырка $b=0$
и кипящий фон $z$, аксиома~A5).
Планкон --- возбуждение $b=1$ \emph{на} этом океане, а не в «мёртвом'' вакууме
с $\Phi\equiv 0$.

\paragraph{Итерация закона $g$.}
Точка $\gamma_t=(q_t,p_t)\in\Gamma_{\hv}$; за один такт на~$\layerMicro$ закона $g$ получаем $\gamma_{t+1}$
(не путать закон~$g$ с узлом, где сидит планкон).
На кипящем океане \emph{внутреннее состояние} планкона меняется по $\Gamma_{\hv}$: различные $(q,p)$
и ненулевой импульс $\Phi$ за десятки тактов, при сохранении топологии $n=\pm1$ в срезе наблюдения.

Орбиты одного планкона на одном узле (основная и возбуждённые ветки по $n_E$) согласованы с эволюцией~$g$.
Период \emph{всего поля} на решётке~$\Lambda$ --- отдельная задача (\cref{sec:floor0-open,para:floor0-two-periods}).

\section{Статус этажа~0}
\label{sec:floor0-open}

Критерий этажа~0: алгебра вех (\cref{sec:floor0-ne-ladder,prop:floor0-ne-ladder}),
основное состояние (\cref{prop:floor0-vortex-ground}), монодромия $SU(2)$
(\cref{prop:floor0-su2-monodromy}), каталог регистров (\cref{sec:floor0-catalog})
и внутренняя динамика \emph{одного} планкона на \emph{одном} узле на кипящем океане
(\cref{sec:floor0-phase-space}).

\paragraph{Закрыто.}
После релаксации вихря на ядре остаётся $n_E=0$ (\cref{prop:floor0-vortex-ground});
при дальнейшей эволюции закона~$g$ на кипящем океане (A5) на том же узле
появляются такты с $n_E\ge 1$, орбиты в $\Gamma_{\hv}$ разделяются на основную и возбуждённые ветки,
а правила отбора между уровнями $n_E$ из \cref{ch:evolution} не нарушаются.

\begin{table}[htbp]
  \centering
  \caption{Закрытые вопросы этажа~0 (один планкон, один $\hv$, кипящий океан).}
  \label{tab:floor0-closed}
  \begin{tabular}{@{}p{0.38\linewidth}p{0.52\linewidth}@{}}
    \toprule
    вопрос & итог \\
    \midrule
    лестница $n_E$, вехи $41,82,256,512$ &
    \cref{prop:floor0-ne-ladder}, \cref{tab:floor0-landmarks} \\
    основное состояние после релаксации &
    $n_E=0$ при $b=1$, $|n|$ на контуре (\cref{prop:floor0-vortex-ground}) \\
    возбуждения $n_E\ge 1$ при свободной эволюции &
    после релаксации в ходе эволюции; не в основном срезе \\
    орбиты $\Gamma_{\hv}$, основная и возбуждённые ветки &
    десятки различимых $(q,p)$ при сохранении $n=\pm1$ \\
    отбор между уровнями $n_E$ &
    правила согласованности \cref{ch:evolution}; проверено на ядре \\
    \bottomrule
  \end{tabular}
\end{table}

\paragraph{Открыто.}

\paragraph{Два разных «периода».}
\label{para:floor0-two-periods}
Число $512$ на этаже~0 --- это размер фазового кольца $\mathbb{Z}_{\phaseRingCount}$
и длина \emph{полного оборота} импульса~$\Phi$ в \emph{одной} планковской ячейке~$\hv$
(\cref{sec:floor0-ne-ladder}).
При постоянном шаге $\Delta\varphi_{\mathrm{disc}}=41$ по этому кольцу тот же вычет~$\Phi$
возвращается через $512$ M-тиков: это \emph{арифметика внутреннего регистра} планкона,
а не свойство всей решётки.
Иное --- \emph{период автомата на всей}~$\Lambda$:
через сколько тактов одновременно на \emph{всех} узлах тора повторяется пара состояний
схемы «лягушка'' при том же законе~$g$ на кипящем океане.
Голономия и projected collision связывают соседей, поэтому оборот кольца на одном~$\hv$
не влечёт повтор поля и не задаёт глобальный цикл.
Лестница $n_E$ и вехи $41,82,256,512$ относятся к первому смыслу;
второй остаётся открытым и не входит в критерий закрытия этажа~0.

\begin{table}[htbp]
  \centering
  \caption{Открытые вопросы (этаж~0 и соседние этажи).}
  \label{tab:floor0-open}
  \begin{tabular}{@{}p{0.38\linewidth}p{0.52\linewidth}@{}}
    \toprule
    вопрос & статус \\
    \midrule
    период автомата на всей решётке $\Lambda$ &
    открыт: точный цикл пары состояний на всех узлах и период макроскопических наблюдаемых
    на аттракторе (\cref{para:floor0-two-periods}); не тождественен обороту $\Phi$ на одном~$\hv$ \\
    полный перечень $2^{B_{\hv}}$ классов &
    не требуется; достаточно орбит на одном узле \\
    внешние моды на $\varepsilon$-оболочке (1-я сфера соседей) &
    \cref{ch:floor1} \\
    \bottomrule
  \end{tabular}
\end{table}

\paragraph{Связь с электроном.}
Лабораторный электрон --- частица на~$\layerMacro$; планкон --- носитель на~$\layerMicro$
(\cref{para:floors-three-carriers,para:floors-composite-ladder}).
Масса, заряд и радиус (через облако $\rho_\Theta$, \cref{ch:matter}) выводятся на других этажах.
Внутренний спектр этажа~0 --- \emph{микроскопический прототип зонной структуры}
(дискретные уровни, отборы, Паули), общий для любого планкона,
в том числе для того, что на~$\layerMacro$ читается как электрон.

\paragraph{Связь с кварками и квантовыми числами.}
Цвет, поколение, изоспин --- не здесь.
Их разбор --- после закрытия этажа~0 и внешних оболочек (этаж~1),
когда составной узел на~$\varepsilon$ даёт больше степеней свободы, чем одно~$\hv$.
